%%=============================================================================
%% Gewenste situatie (TO-BE)
%%=============================================================================

\chapter{\IfLanguageName{dutch}{Gewenste situatie}{Target situation}}%
\label{ch:gewenste-situatie}

Op basis van de tekortkomingen die in het vorige hoofdstuk werden geïdentificeerd, beschrijft dit hoofdstuk de vereisten voor de oplossing en de architectuurvisie voor het gewenste systeem.

%%----------------------------------------------------------------------
\section{Vereisten voor de oplossing}%
\label{sec:tobe-vereisten}

\subsection{Functionele vereisten}

% TODO: Lijst de functionele vereisten op:
% - Automatisch aanmaken van een containerinstantie per gebruiker/team bij
%   het starten van een challenge
% - Automatisch stoppen en verwijderen na afloop of timeout
% - Configureerbare resource-limieten per challenge (CPU, RAM)
% - Netwerkisolatie tussen instanties van verschillende gebruikers
% - Ondersteuning voor meerdere challenge-types
% - Gebruiksvriendelijke interface voor organisatoren om challenges toe te voegen

\subsection{Niet-functionele vereisten}

% TODO: Lijst de niet-functionele vereisten op:
% - Volledig self-hosted en open-source
% - Geen Kubernetes als primaire vereiste
% - Schaalbaar over meerdere VM's
% - Veilig: minimale aanvalsoppervlakte, container escapes voorkomen
% - Beheerbaar: monitoring en logging van actieve instanties
% - Reproduceerbaar: gedocumenteerde installatie-instructies

%%----------------------------------------------------------------------
\section{Architectuurvisie}%
\label{sec:tobe-architectuur}

% TODO: Beschrijf de architectuur van het gewenste systeem:
% - CTFd communiceert via plugin met Docker Engine API
%   (\autocite{DockerEngineAPI}, \autocite{DockerSDKPython})
% - Containers worden automatisch aangemaakt/vernietigd op basis van
%   gebruikersacties in CTFd
% - Resource-limieten worden afgedwongen bij aanmaak
% - Netwerkisolatie via custom Docker bridge networks
% - Reverse proxy (Traefik \autocite{Traefik} of Nginx) als entry point
% - Systeem is modulair: nieuwe challenge-types toevoegen zonder kerncode
%   te wijzigen

% TODO: Voeg hier een componentendiagram in.

%%----------------------------------------------------------------------
\section{Granulariteitsbeslissing: container per gebruiker of per team}%
\label{sec:tobe-granulariteit}

% TODO: Bespreek de voor- en nadelen van beide opties:
% - Per gebruiker: maximale isolatie, meer resource-gebruik
% - Per team: minder instanties, risico op interne interferentie
% Verantwoord de keuze die in de PoC wordt geïmplementeerd.

%%----------------------------------------------------------------------
\section{Schaalbaarheidsvereisten}%
\label{sec:tobe-schaalbaarheid}

% TODO: Beschrijf wat het systeem moet aankunnen:
% - Hoeveel gelijktijdige deelnemers?
% - Hoeveel actieve containers tegelijk?
% - Hoe worden containers verdeeld over meerdere hosts?
% - Docker Swarm \autocite{DockerSwarm} of andere orchestratielaag?
